% STEG-Results.tex -- results for the Rwanda tourism revenue sharing project.
% Figures are the note-free copies written by output/figures/make_tex_figures.py into
% <Dropbox>/output/figures/tex/. Each figure carries its caption above and its note below, here in
% the document rather than inside the image.
\documentclass[11pt]{article}
\usepackage[margin=1in]{geometry}
\usepackage{graphicx,amsmath,amssymb,booktabs,caption,float,microtype}
\usepackage[hidelinks]{hyperref}
\captionsetup[figure]{labelfont=bf,textfont=normalfont,justification=raggedright,singlelinecheck=false,skip=6pt}
\graphicspath{{/Users/matteo/Library/CloudStorage/Dropbox/1-Ongoing Projects/Rwanda - TRS/output/figures/tex/}}
\newcommand{\fignote}[1]{\vspace{4pt}\begin{minipage}{\linewidth}\footnotesize\setlength{\parskip}{0pt}\textit{Notes.} #1\end{minipage}}
\setlength{\parskip}{5pt}
\title{Tourism revenue sharing around Rwanda's national parks\\[2pt]\large Results}
\author{}
\date{\today}

\begin{document}
\maketitle

\section{What the design is trying to separate}

Rwanda's tourism revenue sharing scheme returns a share of national park receipts to the communities
that border the parks. Every bordering sector receives the transfer. Only a few of them also sit at a
park entrance, where tourists actually arrive, sleep and spend. That distinction is the whole
identification strategy of this project: it lets the money be separated from the tourism.

Three groups, defined at sector level because that is the level the transfer is allocated at:

\begin{center}
\begin{tabular}{llr l}
\toprule
& Definition & $n$ & Channel \\
\midrule
$G_1$ & Gates+: borders a park \emph{and} sits within 5\,km of an entrance & 12 & Transfer \emph{and} tourism \\
$G_2$ & The other bordering sectors & 32 & Transfer only \\
$G_0$ & Neither & 319 & Comparison group \\
\bottomrule
\end{tabular}
\end{center}

$G_1$ against $G_0$ measures the two channels together, $G_2$ against $G_0$ measures the transfer on
its own, and $G_1-G_2$ is what tourism adds. The three groups plus the exclusions below account for
all 416 sectors: $12+32+319+9+44=416$.

\paragraph{Exclusions, and why.} The City of Kigali and the four largest towns of 2002 (44 sectors)
are dropped: they are urban outliers on every outcome and none of them borders a park, so they add
nothing to the comparison but a great deal of variance. The nine sectors bordering Gishwati-Mukura
are also dropped. That park was gazetted only in 2015--16, so those sectors enter the revenue-sharing
zone a decade after the others; leaving them among the controls would place nine treated units in the
comparison group precisely in the years where the estimates move.

\paragraph{A caution carried through the whole document.} $G_1$ contains twelve sectors and is
identified off five districts. Standard errors clustered at sector assume many more independent
groups than that, so the $p$-values attached to $G_1$ should be read as optimistic. A wild cluster
bootstrap is the right fix and has not yet been run.

\section{Housing quality}

The housing outcomes use the census in an ANCOVA, conditioning on the 2002 level of the same outcome
rather than differencing it:
\begin{equation}
Y_{s,t} \;=\; \alpha \;+\; \beta\,\text{Treated}_s \;+\; \gamma\,Y_{s,2002} \;+\; \phi_{p(s)} \;+\; \delta_{d(s)} \;+\; \varepsilon_{s,t},
\qquad t \in \{2012,\,2022\}
\end{equation}
where $\phi_{p(s)}$ is a nearest-park fixed effect and $\delta_{d(s)}$ a district fixed effect; the
pooled specification replaces the latter with district-by-wave effects. Errors are clustered at
sector and reported alongside Conley spatial standard errors at 25 and 50\,km. Conditioning on the
2002 level absorbs the pre-treatment differences that a bare cross-section would attribute to
treatment, which matters here because park-adjacent sectors were poorer than average in 2002.

\begin{figure}[H]
\caption{Housing components and the housing index, by treatment definition}
\includegraphics[width=\linewidth]{bar_chart_housing.pdf}
\fignote{Bars are the change between census waves in the share of households with each housing
attribute, for each treatment definition against the control sectors. The index is the unweighted
mean of the five standardised components. Confidence intervals are from a sector-level bootstrap.
The City of Kigali and the four largest towns of 2002 are excluded.}
\end{figure}

\begin{figure}[H]
\caption{Housing: ANCOVA estimates conditioning on the 2002 level}
\includegraphics[width=\linewidth]{coefplot_housing.pdf}
\fignote{Each point is $\beta$ from equation~(1), in units of the 2002 control standard deviation,
with 95\% confidence intervals from sector-clustered standard errors. The specification includes the
2002 level of the outcome, nearest-park fixed effects and district fixed effects; the pooled column
uses district-by-wave fixed effects. Modern roofing is the outcome that moves: $+0.366$ standard
deviations pooled for the 44 bordering sectors ($p<10^{-7}$), and the housing index rises $+0.277$
($p=0.016$). Water and sanitation do not move.}
\end{figure}

\begin{figure}[H]
\caption{Housing: the same estimates with controls chosen by post-double-selection lasso}
\includegraphics[width=\linewidth]{coefplot_housing_lasso.pdf}
\fignote{As the previous figure, with the control set chosen by post-double-selection lasso with
cluster-robust penalty loadings, drawn from the 2002 sector characteristics. The lasso is run
separately for each outcome and wave; the number of selected controls is reported in the underlying
table. The roofing result survives control selection, which is the point of running it.}
\end{figure}

\section{Employment}

Employment uses the same ANCOVA, with the outcome the inverse hyperbolic sine of the count of workers
aged 20--64 in each of the three Fisher--Clark sectors, so coefficients read as proportional changes.

\paragraph{The classification problem, and what was done about it.} The three censuses do not use the
same industry coding: 2002 is ISIC rev.\ 3 at the group level and 2012 and 2022 are ISIC rev.\ 4
sections. The mapping was rebuilt so that the three sectors mean the same thing in all three waves,
and verified to be exhaustive and non-overlapping in each. Mining is placed in the primary sector.

\paragraph{The 2022 break, and the reconstruction.} The 19th ICLS revision, applied in 2022, counts
agricultural self-employment only when production is \emph{mainly for market}. Subsistence farmers
therefore leave the recorded workforce entirely in 2022, and recorded primary employment falls from
71.6\% to 60.1\% to 27.3\% of adults across the three waves --- a break in measurement, not in
behaviour. The 2022 primary count is reconstructed by adding back adults with no recorded industry
who live in a crop-growing household, less the 14.1 percentage point false-positive rate that the
same rule produces in 2012, where subsistence farmers \emph{are} recorded. The reconstruction gives a
monotone 71.6 $\rightarrow$ 60.1 $\rightarrow$ 46.1\%, and is validated by the fact that the
farm-household composition of unrecorded adults is 67.5\% in 2012 against 67.8\% in 2022.

\begin{figure}[H]
\caption{Employment by broad sector: ANCOVA estimates}
\includegraphics[width=\linewidth]{coefplot_jobs.pdf}
\fignote{Outcomes are $\operatorname{asinh}$ of the count of workers aged 20--64, so coefficients are
approximately proportional changes, expressed in 2002 control standard deviations. Specification as
in equation~(1). For the 44 bordering sectors the pooled estimates are $+0.576$ for primary
($p=0.0001$), $+0.160$ for secondary ($p=0.016$) and $+0.177$ for tertiary ($p=0.042$): employment
rises across the board, which is consistent with a labour market expanding rather than reallocating.
The 2022 primary count is the reconstructed series described above.}
\end{figure}

\begin{figure}[H]
\caption{Employment by broad sector, controls chosen by post-double-selection lasso}
\includegraphics[width=\linewidth]{coefplot_jobs_lasso.pdf}
\fignote{As the previous figure with lasso-selected controls. The Gates+ pooled estimates change sign
on primary employment once controls are selected: primary $-0.602$ ($p=0.018$), secondary $+0.184$
($p=0.056$), tertiary $+0.265$ ($p=0.050$). Read together with the unselected estimates, the pattern
at the gates is a shift out of agriculture and into services rather than a uniform expansion; the
sign flip on primary is driven by control selection and should be reported as such rather than as a
robust result.}
\end{figure}

\section{Education}

Schooling is a stock fixed by the time a person finishes school, so a difference-in-differences on
the adult stock would attribute to treatment things that happened decades earlier. The design instead
follows the cohort-exposure logic of Duflo (2001): a person's exposure depends on how old they were
when revenue sharing began. Children who had already left school by 2005 are unexposed; those who
started primary school in 2005 or later are fully exposed.
\begin{equation}
Y_{i,c,s} \;=\; \alpha \;+\; \sum_{b}\beta_b\,\bigl[\text{Treated}_s \times \mathbf{1}(c \in b)\bigr]
\;+\; \delta_c \;+\; \mu_s \;+\; X_{i}'\theta \;+\; \varepsilon_{i,c,s}
\end{equation}
with $c$ the birth cohort, $\delta_c$ cohort fixed effects, $\mu_s$ sector fixed effects, and $b$
indexing three-year cohort bins. The reference cohort is 1999, the first to start primary school in
2005. Cohorts before 1970 are dropped. Because several cohorts were of school age during the
genocide, the specification is also run controlling for sector-level violence exposure --- double
orphanhood, any orphanhood, war disability, widowhood, the sex ratio, child-headed households, prior
residence abroad and foreign birth --- entered individually rather than as an index.

\begin{figure}[H]
\caption{Education by birth cohort: years of schooling, completed primary, and literacy}
\includegraphics[width=\linewidth]{event_study_education_3yr.pdf}
\fignote{Each point is the coefficient on treatment interacted with a three-year birth-cohort bin
from equation~(2), relative to the 1999 cohort, with 95\% intervals from sector-clustered standard
errors. Cohorts to the left of the reference were too old to be affected and act as placebos. At the
gates the pooled cohort estimates are $+0.337$ years of schooling ($p=0.0014$) and $+2.4$ percentage
points on completed primary ($p=0.023$); the 44-sector definition is null on all three outcomes. The
violence controls reduce the point estimates modestly and do not overturn them.}
\end{figure}

\section{Age at marriage}

Marriage uses the same cohort design, restricted to women aged 18 and over, with cohorts running to
2004. The outcomes are an indicator for having married by 18 and age at first marriage.

\begin{figure}[H]
\caption{Marriage by birth cohort: married by 18 and age at first marriage}
\includegraphics[width=\linewidth]{event_study_marriage_3yr.pdf}
\fignote{Specification as in equation~(2), three-year cohort bins, reference cohort 1999, 95\%
intervals from sector-clustered standard errors. At the gates, marriage by 18 falls $3.6$ percentage
points ($p=0.0012$) and age at first marriage rises $0.42$ years ($p=0.004$); for the 44 bordering
sectors the corresponding figures are $-2.6$ points ($p=0.0007$) and $+0.33$ years ($p=0.0003$).
Controlling for genocide violence exposure leaves both intact ($-3.3$ points, $p=0.0010$; $+0.34$
years, $p=0.008$). Pre-trends are flat. This is the best-identified result in the project.}
\end{figure}

\section{Forest loss}

Forest uses a fixed 1\,km grid rather than administrative units, and a matched design rather than a
bare difference-in-differences, because treated and control cells differ systematically in terrain,
accessibility and baseline forest stock:
\begin{equation}
L_{i,t} \;=\; \alpha_i \;+\; \delta_{d(i),t} \;+\; \sum_{t}\beta_t\bigl[\text{Treated}_i \times \mathbf{1}(\text{year}=t)\bigr]
\;+\; W_{i,t}'\gamma \;+\; \varepsilon_{i,t}
\end{equation}
estimated with overlap weights and weighted by them. $L_{i,t}$ is hectares of annual forest loss,
$\alpha_i$ a cell fixed effect, $\delta_{d(i),t}$ district-by-year fixed effects, and $W_{i,t}$
annual rainfall, temperature and PDSI. Overlap weights are estimated separately for each panel from
2004 forest stock, elevation, slope, distance to trunk or primary roads, 2001 gridded population
density, baseline precipitation, and the 2001--2004 level and trend of forest loss. Standard errors
are clustered on 20-by-20\,km spatial blocks. Cells must hold at least 10 hectares of the relevant
forest measure in 2004. The reference year is 2004 and the two products --- Hansen Global Forest
Change and JRC Tropical Moist Forest --- are estimated separately, since they disagree about what
counts as forest.

\begin{figure}[H]
\caption{Forest loss around the parks: annual event-study estimates, five comparisons}
\includegraphics[width=\linewidth]{event_study_forest_5km.pdf}
\fignote{Points are year-specific treated-control differences in annual forest loss relative to 2004,
from equation~(3); bars are 90\% intervals from 20-by-20\,km spatial-block clustered standard errors.
The first three panels are geographic; the last two are the sector split of Section~1 and compare
each group against the same far cells, less the sectors bordering Gishwati-Mukura. Protection works
inside the boundary --- pooled post-2005, Hansen gives $-0.365$ hectares per cell ($p<0.001$) and JRC
$-0.275$ ($p=0.006$) for park against non-park --- but neither the transfer nor tourism shows outside
it: Gates+ sectors give $+0.027$ ($p=0.50$, Hansen) and $-0.035$ ($p=0.90$, JRC), and
revenue-sharing-only sectors $-0.030$ ($p=0.73$) and $+0.019$ ($p=0.94$). Pre-trend $p$-values
jointly test the 2001--2003 coefficients.}
\end{figure}

\begin{figure}[H]
\caption{Forest loss: pooled post-2005 estimates from the matched-pair design}
\includegraphics[width=\linewidth]{coefplot_forest_5km.pdf}
\fignote{Pooled post-2005 coefficients from a matched-pair difference-in-differences in which each
treated cell is matched with replacement to the closest eligible control cell in the same district on
the covariates listed above, with pair fixed effects and baseline-covariate-specific linear trends.
Bars are 90\% intervals with two-way clustering on treated and matched-control spatial blocks.
Negative coefficients mean less forest loss in the first area. This figure covers the geographic
comparisons only: its coefficients are stored constants from an estimator whose code is not in the
repository, so it could not be regenerated for the sector split.}
\end{figure}

\begin{figure}[H]
\caption{Forest loss measured by JRC Tropical Moist Forest}
\includegraphics[width=\linewidth]{coefplot_jrc_forest.pdf}
\fignote{The JRC product measured on its own terms, as a check that the Hansen results are not an
artefact of how Hansen defines tree cover. JRC counts moist forest and excludes plantation, which is
why its levels differ from Hansen even where the two agree in sign.}
\end{figure}

\section{Nighttime lights}

Lights are the one outcome available annually for the whole period, which makes them the natural
place to look for a time path. Both groups enter one equation so that they share a comparison group:
\begin{equation}
Y_{u,t} \;=\; \alpha_u \;+\; \delta_{d(u),t}
\;+\; \sum_{t}\beta^{1}_{t}\bigl[G_1(u)\times\mathbf{1}(\text{year}=t)\bigr]
\;+\; \sum_{t}\beta^{2}_{t}\bigl[G_2(u)\times\mathbf{1}(\text{year}=t)\bigr] \;+\; \varepsilon_{u,t}
\end{equation}
$Y_{u,t}$ is $\operatorname{asinh}$ of the sector's total lights, scaled by the 2002 control standard
deviation. $\alpha_u$ absorbs everything fixed about a place and $\delta_{d(u),t}$ every district-level
shock, so the comparison is between sectors in the same district in the same year. The reference year
is 2004 and errors are clustered at sector. The tourism increment $\beta^1-\beta^2$ is estimated
directly, with its own standard error, by re-running the same equation on $[G_1 \text{ or } G_2]$ and
$[G_1]$.

\paragraph{Five restrictions, each settled by a test rather than assumed.}
\begin{enumerate}\itemsep2pt
\item \textbf{Park area is removed from every unit before the lights are counted.} Park land is dark
by construction and it is not evenly spread: treated sectors lose 31.6\% of their area to park
against 0.2\% for controls, so leaving it in attenuates the treated units specifically.
\item \textbf{DMSP is put on one satellite's scale before the satellites are averaged.} Twelve years
have two satellites flying, and averaging them raw treats two instruments with different gain as if
they agreed. A monotone $DN\rightarrow DN$ map is fitted per satellite by quantile matching on
Rwanda's own pixels in the months when both saw the same ground, and chained onto F15. F18
(2010--2013) never overlaps another DMSP satellite, so it cannot be chained by measurement and stays
raw; the seam at 2010 is left visible rather than papered over.
\item \textbf{Raw sensors only; no harmonised product.} DMSP-OLS ended in 2013 and VIIRS begins in
2012, so the two are estimated separately and VIIRS is shifted by the DMSP coefficient at 2012, its
intervals carrying that offset's variance. The splice is licensed by the 2012--13 overlap, where both
sensors observe the same units and the treated-control gap between them is $-0.018$ ($p=0.89$). The
Li and Chen harmonised products are \emph{not} used: each is half model output, and their large
post-2013 effects live entirely in the modelled half.
\item \textbf{Uncensored \texttt{avg\_vis} rather than \texttt{stable\_lights}}, which zeroes dim
cells --- and the treated units are the dark ones.
\item \textbf{Sector is the headline unit} because treatment is assigned at sector level; the cell
panel is shown beside it to confirm the weighting does not drive the result, not as extra
identification.
\end{enumerate}

\paragraph{Why the $p$-values here come from randomisation.} $G_1$ is twelve sectors in five
districts and the gate design below is thirty-three cells in five, all strung along park borders.
Clustering at sector assumes those units are independent draws within a year; they are not, since a
good year at one park border is a good year at all of them. The consequence is measurable rather than
hypothetical: reassigning the group label at random among the controls and recomputing the joint
pre-trend statistic, random twelve-sector groups produce a \emph{median} Wald of 35 to 80 against a
$\chi^2$ reference with twelve degrees of freedom. Every headline coefficient and every pre-trend test
below therefore carries a randomisation $p$-value, and that is the one to read. On this correction
the pre-trends pass comfortably for $G_1$ ($p=0.87$ annual, $0.84$ monthly) and sit at the margin for
$G_2$ ($p=0.06$), where the uncorrected test had reported $p=0.0006$.

\paragraph{A measurement limit that bounds what lights can show.} Across the twelve years when two
DMSP satellites flew simultaneously, the same sectors are observed twice. The two readings disagree
about the treated-control gap by \emph{more} than that gap varies across years: test-retest
correlation is 0.45 at best and about 0.10 for the twelve gate sectors. Intercalibration raises those
to 0.33 and 0.23 without rescuing them. The DMSP era therefore cannot support a small effect, and its
flatness is weak evidence of a null rather than strong evidence of one.

\begin{figure}[H]
\caption{Nighttime lights: revenue sharing with tourism against revenue sharing alone}
\includegraphics[width=\linewidth]{event_ntl_groups.pdf}
\fignote{Year-specific coefficients from equation~(4), relative to 2004, in 2002 control standard
deviations, with 95\% intervals from sector-clustered standard errors. Blue is DMSP-OLS 1992--2013 on
the intercalibrated F15 scale and amber VIIRS 2012--2025, spliced at the dashed line. Post-2005
differences-in-differences, estimable on DMSP only because VIIRS begins seven years after treatment:
$G_1$ $+0.009$ ($p=0.78$), $G_2$ $-0.023$ ($p=0.14$), and the tourism increment $G_1-G_2$ $+0.032$
($p=0.34$). Nothing is distinguishable from zero on either channel. The visible movement is 2019
onward, fourteen years after revenue sharing began, and is common to both groups.}
\end{figure}

\subsection*{The gates, at 5 km}

Sectors are the wrong unit for a tourism effect. They average 65 km$^2$ of settled land, and the
lodges and service clusters a gate would generate occupy a small fraction of that, so a sector total
dilutes whatever is there. The same design is therefore run on cells, comparing every cell within
5 km of a park gate against every cell beyond it, and separately every cell within 5 km of the park
itself:
\begin{equation}
Y_{c,t} \;=\; \alpha_c \;+\; \delta_{d(c),t}
\;+\; \sum_{t}\beta_{t}\bigl[\text{within5}(c)\times\mathbf{1}(\text{year}=t)\bigr] \;+\; \varepsilon_{c,t}
\end{equation}
Cell fixed effects and district-by-year fixed effects. Sector-by-year effects are \emph{not} used:
sectors are smaller than the scale this treatment varies over, so sector-by-time absorbs the
treatment itself --- only 6 of 33 gate cells keep identifying variation under it, against 33 of 33
under district-by-year.

\begin{figure}[H]
\caption{Nighttime lights within 5 km of a park gate, and within 5 km of the park}
\includegraphics[width=\linewidth]{event_ntl_gates.pdf}
\fignote{Year-specific coefficients from equation~(5), relative to 2004, in 2002 control standard
deviations, with 95\% intervals. DMSP-OLS on the intercalibrated F15 scale to 2013, VIIRS thereafter,
spliced at the dashed line. The post-2005 difference-in-differences is $+0.033$ at the gate, with a
clustered $p$ of $0.089$ and a randomisation $p$ of $0.247$: \emph{not} a result. An earlier version
of this design reported $p=0.0496$; that figure rested on raw satellite scales and on standard errors
that assume thirty-three contiguous cells in five districts are independent. What does survive is
growth in the VIIRS era: gate cells brighten $+0.147$ control SD per year faster than cells beyond
5 km over 2012--2025 (randomisation $p=0.005$), and $+0.124$ over 2017--2025 ($p=0.005$), the window
that excludes VIIRS's noisy early years. Proximity to the \emph{park} shows nothing on any of these
($+0.011$, $p=0.58$; $+0.013$, $p=0.36$) despite resting on 227 treated cells rather than 33.}
\end{figure}

Two checks decide how much weight this carries. The growth specification sits entirely after
treatment, so it cannot test parallel trends on its own; the monthly panel, which has about 150
pre-2005 observations per cell, can. Gate cells show \textbf{no} differential pre-trend
(randomisation $p=0.84$), so the post-2012 divergence is not the continuation of a path they were
already on. Park-proximity cells \textbf{fail} the same test ($p=0.005$). The gate/park contrast
should therefore be read as \emph{the gate design is identified and the park design is not}, rather
than as a like-for-like discriminant between two valid comparisons.

\subsection*{Monthly composites}

The annual design rests on thirteen pre-treatment observations per unit. The monthly composites give
about 150, which is what makes a pre-trend test informative rather than decorative. Nothing else
changes: same groups, same settled land, same reference year, same clustering, with district-by-%
\emph{month} fixed effects so that the seasonal cycle is absorbed --- Rwanda's cloud-free night count
runs from 2.5 in April to 8.8 in July, and seasonality accounts for 40\% of residual variance in
DMSP.

\begin{figure}[H]
\caption{Nighttime lights by month, 1992--2025: revenue sharing with tourism against revenue sharing alone}
\includegraphics[width=\linewidth]{event_ntl_monthly_prepost.pdf}
\fignote{Year-specific coefficients estimated on 416 sectors by 401 months of settled land, relative
to 2004, in control standard deviations, with 95\% intervals. Totals are scaled to each sector's whole
settled area: summing only observed pixels would measure coverage as much as light, and treated
sectors are observed less (79\% of sector-months clear a 90\% threshold against 86\% for controls), a
shortfall that halved between 1992--94 and 2005--09. Pre/post estimates on the calibrated F10--F16
block, which contains the entire pre-period and never crosses the 2010 seam: $G_1$ $+0.030$
($p=0.32$, randomisation $p=0.20$), $G_2$ $+0.007$ ($p=0.66$), increment $+0.024$ ($p=0.45$).
Pre-trends by randomisation: $G_1$ $p=0.84$, $G_2$ $p=0.055$. Unlike the annual case, the monthly
DMSP/VIIRS join is \emph{not} licensed by the overlap: across the twenty F18/VIIRS overlap months the
two sensors disagree about the treated-control gap (correlation $-0.26$ and $-0.11$; mean difference
$-0.44$ and $-0.37$ control SDs, $p=0.005$ and $0.013$). The join drawn here is a plotting convention
and no coefficient crosses it.}
\end{figure}

\begin{figure}[H]
\caption{The 5 km gate design on monthly composites}
\includegraphics[width=\linewidth]{event_ntl_gates_monthly.pdf}
\fignote{Equation~(5) on 1{,}939 cells by 401 months, cell and district-by-month fixed effects.
Magnitudes are not comparable with the annual gate figure: each series is scaled by its own control
standard deviation, and the monthly composites are far more dispersed than the cleaned annual
product. Sign and significance are what transfer. Gate growth over 2017--2025 replicates
(randomisation $p=0.010$ against $0.005$ annually) while park proximity stays null ($p=0.45$ against
$0.36$). Over 2012--2025 the monthly panel instead shows a park effect the annual data does not,
which is the window that includes VIIRS's noisy early years: residual dispersion falls fivefold at
2017, from 0.96 in 2012--14 to 0.18 in 2017--18.}
\end{figure}

\section{What the results add up to}

The transfer arm and the tourism arm behave differently, and not in the way the programme's own
theory of change would predict.

\begin{enumerate}\itemsep3pt
\item \textbf{The best-identified results are demographic, not economic.} Marriage by 18 falls and
age at first marriage rises at the gates, on flat pre-trends, robust to genocide-violence controls.
Years of schooling rise at the gates and nowhere else. Both come from the cohort design, which has
far more identifying variation than the three-wave census panels.
\item \textbf{Housing improves, but narrowly.} Roofing moves and survives control selection; water and
sanitation do not move at all. That is the signature of household-level investment rather than public
infrastructure, which is what a cash transfer to communities would be expected to produce only if it
were reaching households.
\item \textbf{Employment expands rather than reallocating}, except at the gates, where the lasso
estimates point to a shift out of agriculture into services. The sign flip on primary employment
under control selection is not robust and should not be reported as a finding on its own.
\item \textbf{Forest responds to neither channel; lights respond only at the gates, and only
recently.} Protection works inside the park boundary on both products; outside it, neither the
transfer nor tourism moves forest loss. Lights are flat on both arms at sector level, in the annual
and the monthly design alike, with the caveat above that the DMSP era cannot detect a small effect.
The exception is spatial and late: cells within 5 km of a gate brighten $+0.147$ control SD per year
faster than the rest of the country over 2012--2025 (randomisation $p=0.005$), on flat pre-trends,
while proximity to the park does nothing. That the result appears only at cell level is what dilution
predicts --- thirty-three cells inside sectors averaging 65 km$^2$ --- and it is the one lights result
that survives correcting the inference.
\item \textbf{Tourism adds little that the transfer does not, except at the gate itself.} Where
$G_1$ and $G_2$ are compared directly at sector level --- lights and forest --- the increment is
indistinguishable from zero. Where the gates stand out --- schooling, marriage, the services shift,
and now lights within 5 km --- the effect appears in outcomes a transfer alone would not obviously
move, and it is concentrated in the few kilometres where visitors actually arrive. That points at
labour-market opportunity rather than money, and at a catchment far smaller than the sector the
revenue-sharing formula pays.
\end{enumerate}

\paragraph{What would change these conclusions.} A wild cluster bootstrap on every $G_1$ estimate,
since twelve treated sectors in five districts cannot support cluster-robust inference as reported.
An explanation for the 2019--2022 break in lights, which appears in every specification and every
product and currently sits underneath the only positive lights result. And a measure of rural
electrification by sector: nighttime lights in a poor rural country are substantially a measure of
where the grid went, and without that control a lights result cannot be attributed to this programme.

\end{document}
